\documentclass[12pt]{article}
\usepackage{../../ficha}

\begin{document}

\ficha{Sistemas de Equações}[Sistemas de Duas Equações do 1.º Grau a Duas Incógnitas]

\instrucoes{Resolve os sistemas, apresentando todos os passos. Verifica as soluções sempre que possível.}

\secao{Verificar Soluções}{Verifica se o par ordenado indicado é solução do sistema.}
\begin{exercicios}
  \item $(2, 1)$ \; em \; $\begin{cases} x + y = 3 \\ x - y = 1 \end{cases}$
  \item $(3, -1)$ \; em \; $\begin{cases} 2x + y = 5 \\ x + 3y = 1 \end{cases}$
  \item $(1, 4)$ \; em \; $\begin{cases} y = 3x + 1 \\ x + y = 6 \end{cases}$
  \item $(-2, 3)$ \; em \; $\begin{cases} 3x + 2y = 0 \\ x - y = -5 \end{cases}$
  \item Indica três soluções da equação $x + 2y = 8$.
\end{exercicios}

\secao{Método de Substituição}{Resolve pelo método de substituição.}
\begin{exercicios}
  \item $\begin{cases} y = 2x \\ x + y = 9 \end{cases}$
  \item $\begin{cases} x = y + 3 \\ x + y = 11 \end{cases}$
  \item $\begin{cases} y = 3x - 1 \\ 2x + y = 9 \end{cases}$
  \item $\begin{cases} x + y = 7 \\ 2x - y = 2 \end{cases}$
  \item $\begin{cases} x - 2y = 1 \\ 3x + y = 10 \end{cases}$
  \item $\begin{cases} 2x + 3y = 12 \\ x - y = 1 \end{cases}$
  \item $\begin{cases} 4x - y = 5 \\ 3x + 2y = 12 \end{cases}$
  \item $\begin{cases} x + y = -1 \\ 3x - 2y = 12 \end{cases}$
\end{exercicios}

\secao{Sistemas com Parênteses e Frações}{Simplifica as equações e resolve os sistemas.}
\begin{exercicios}
  \item $\begin{cases} 2(x + 1) + y = 10 \\ x - y = 1 \end{cases}$
  \item $\begin{cases} 3(x - y) = 6 \\ x + 2y = 8 \end{cases}$
  \item $\begin{cases} \dfrac{x}{2} + y = 4 \\ x - y = 2 \end{cases}$
  \item $\begin{cases} \dfrac{x}{3} = \dfrac{y}{2} \\ x + y = 10 \end{cases}$
  \item $\begin{cases} \dfrac{x + y}{2} = 5 \\ \dfrac{x - y}{3} = 2 \end{cases}$
  \item $\begin{cases} 2(x - 1) = 3(y + 1) \\ x + y = 10 \end{cases}$
\end{exercicios}

\secao{Classificação de Sistemas}{Classifica cada sistema como possível e determinado, possível e indeterminado ou impossível.}
\begin{exercicios}
  \item $\begin{cases} x + y = 4 \\ x + y = 7 \end{cases}$
  \item $\begin{cases} x + y = 4 \\ 2x + 2y = 8 \end{cases}$
  \item $\begin{cases} x + y = 4 \\ x - y = 2 \end{cases}$
  \item $\begin{cases} y = 2x + 1 \\ 4x - 2y = 3 \end{cases}$
  \item $\begin{cases} 3x - y = 2 \\ -6x + 2y = -4 \end{cases}$
  \item O que se pode dizer das retas que representam um sistema impossível? E um sistema possível e indeterminado?
\end{exercicios}

\secao{Problemas}{Escreve um sistema de equações que traduza o problema e resolve-o.}
\begin{exercicios}
  \item A soma de dois números é 30 e a diferença é 8. Quais são os números?
  \item Num parque há bicicletas e triciclos, num total de 20 veículos e 47 rodas. Quantos há de cada?
  \item Dois cafés e um bolo custam 3,10 euros. Um café e dois bolos custam 3,50 euros. Quanto custa cada um?
  \item Numa quinta há galinhas e coelhos: 35 cabeças e 100 patas. Quantos animais há de cada espécie?
  \item O perímetro de um retângulo é 40\,cm e o comprimento é o triplo da largura. Quais são as dimensões?
  \item Num teatro, os bilhetes de adulto custam 12 euros e os de criança 7 euros. Venderam-se 150 bilhetes por 1\,450 euros. Quantos bilhetes de cada tipo se venderam?
  \item Há 5 anos, a idade do pai era o quádruplo da idade do filho. Hoje, a soma das idades é 55. Que idade tem cada um?
\end{exercicios}

\secao{Desafios}{Resolve os problemas.}
\begin{exercicios}
  \item Resolve: $\begin{cases} \dfrac{x - 1}{2} - \dfrac{y + 2}{3} = 0 \\[6pt] \dfrac{x}{3} + y = 2 \end{cases}$
  \item Determina $k$ de modo que o sistema $\begin{cases} 2x + ky = 3 \\ 4x + 6y = 5 \end{cases}$ seja impossível.
  \item Um número de dois algarismos tem soma dos algarismos 9. Trocando os algarismos, o número aumenta 27. Qual é o número?
  \item Um barco percorre 36\,km a favor da corrente em 2 horas e o mesmo percurso contra a corrente em 3 horas. Qual é a velocidade do barco em águas paradas e a velocidade da corrente?
  \item Uma loja vende dois tipos de café: um a 8 euros/kg e outro a 12 euros/kg. Que quantidade de cada deve misturar para obter 20\,kg a 9,50 euros/kg?
\end{exercicios}

\end{document}
